\section{Related Work} 

\paragraph{Looped transformers \& latent reasoning.} Looped transformers reuse the same set of parameters across multiple forward passes, achieving the computational depth of larger non-looped models without increasing parameter count \citep{dehghani2018universal, zeng2026ponderlm}. Instead of a fixed amount of computation, adaptive looping has also been investigated, where a learned halting mechanism decides how many iterations each input requires \citep{banino2021pondernet, zhu2025scaling, frey2026adaptive}.
Looping in transformers can proceed along two axes both representing a form of latent reasoning. Vertical looping reapplies the same block of layers to iteratively refine the token representations, adding effective depth without lengthening the sequence \citep{saunshi2025reasoning}. 
Horizontal looping instead unfolds across the sequence. A hidden state is fed back as an additional latent token, carrying reasoning into latent space before any discrete token is emitted \citep{hao2024training, shen2025codi}. 
Our approach falls in the vertical category: we loop over the model and aggregate representations across iterations, preserving information rather than overwriting it.








\paragraph{Multi-token prediction.} Next-token prediction provides a single supervision signal per position, which may limit the model's ability to plan ahead \citep{cornille2024learning}. 
Multi-token prediction (MTP) addresses this by training the model to forecast several future tokens at once, yielding denser gradients and encouraging representations that account for longer-range dependencies \citep{liu2024deepseek}.
While MTP has been widely adopted for speculative decoding at inference time, a separate line of work leverages it to improve training itself \citep{gloeckle2024better, liu2024deepseek}. 
With regard to the usage of MTP specifically, the closest work to ours is \citet{noci2026thinking}, who combine MTP with latent reasoning. At selected positions, the model performs a multi-step lookahead in hidden space, supervised against the upcoming ground-truth tokens via a cross-entropy loss over the entire vocabulary.
We reduce this cost by using MTP in latent space as a \textit{soft guiding signal}. 
Rather than predicting a distribution over the vocabulary, we use cosine similarity to align each iteration's hidden state with the embedding of a specific future token. 
Because this signal is applied at every iteration, its cost scales with the number of loops. 
A cross-entropy objective would incur one vocabulary projection per iteration, whereas our cosine alignment adds negligible overhead even at high loop counts. 
Moreover, whereas \citet{noci2026thinking} focus on synthetic planning tasks, we target reasoning in parameter-efficient looped models.
